\documentclass[11pt]{article}

\usepackage[margin=1in]{geometry}
\usepackage{amsmath, amssymb}
\usepackage[T1]{fontenc}
\usepackage[utf8]{inputenc}
\usepackage{lmodern}
\usepackage{enumitem}
\usepackage{hyperref}

\setlist[itemize]{topsep=4pt, itemsep=2pt, parsep=0pt}

\title{6.8610 Spring 2026 Lecture 2 Notes: Log-Linear Language Models and Word Embeddings}
\author{}
\date{February 5, 2026}

\begin{document}
\maketitle

\section{Log-Linear Language Models and Word Embeddings}

\subsection{Motivation}
The lecture introduces the first concrete family of language models used in the course: \textbf{log-linear models}. The core problem is to replace the impossible ``one parameter for every possible context'' view of language modeling with a tractable function approximator.

\begin{itemize}
    \item A language model needs to estimate probabilities of next tokens conditioned on preceding context.
    \item A naive table of probabilities over all possible contexts grows exponentially with sequence length and is unusable.
    \item Log-linear models provide the simplest useful replacement: a softmax over linear scores.
    \item Historically, language models were often used as components inside larger systems such as machine translation, rather than as end-to-end systems on their own.
    \item The lecture then asks how to improve fixed hand-designed context representations, leading to \textbf{word embeddings} and the \textbf{word2vec} family of models.
\end{itemize}

\subsection{Key Concepts}
\textbf{Definition: Log-Linear Language Model}

A log-linear language model predicts the next token by assigning a linear score to each vocabulary item and normalizing with a softmax.

\begin{itemize}
    \item \textbf{Language-model recap:} we want to model
    \[
    p_\theta(x_{1:T}) = \prod_{t=1}^{T} p_\theta(x_t \mid x_{1:t-1})
    \]
    and therefore need a model of \(p_\theta(x_t \mid x_{1:t-1})\).
    \item \textbf{Historical motivation from translation:} older systems often used Bayes' rule to combine a translation model with a language model:
    \[
    p(y \mid x) \propto p(x \mid y)\, p(y)
    \]
    where \(p(y)\) is a language model over target-language strings.
    \item \textbf{What a good language model must know:}
    \begin{itemize}
        \item grammar,
        \item factual knowledge,
        \item common-sense regularities about the world.
    \end{itemize}
    \item \textbf{Prompting view:} once a language model can condition on arbitrary context and continue that context, prompting becomes a general interface for many downstream tasks.
    \item \textbf{Feature representation \(\phi(x)\):} since contexts are variable-length strings, log-linear models require a fixed-dimensional representation of the context before scoring outputs.
    \item \textbf{Softmax / multinomial logistic regression:} the model used here is the multiclass generalization of logistic regression.
    \item \textbf{Embedding models:} instead of using a fixed hand-engineered \(\phi(x)\), we can learn low-dimensional vector representations of words directly from data.
\end{itemize}

\subsection{Core Models and Equations}
\textbf{Autoregressive factorization.}
\[
p_\theta(x_{1:T}) = \prod_{t=1}^{T} p_\theta(x_t \mid x_{1:t-1})
\]

\textbf{Naive scalar parameterization is intractable.} If vocabulary size is \(V\), explicit tables require on the order of
\[
\sum_{t=1}^{T} V^t
\]
parameters.

\textbf{Log-linear next-token model.} Let \(x\) denote the input context and \(y\) the next token. Then
\[
p_\theta(y \mid x) = \frac{\exp\!\left(\theta_y^\top \phi(x)\right)}
{\sum_{y' \in V} \exp\!\left(\theta_{y'}^\top \phi(x)\right)}
\]
where:
\begin{itemize}
    \item \(\phi(x)\) is a fixed-dimensional representation of the context,
    \item \(\theta_y\) is the parameter vector associated with output token \(y\),
    \item \(V\) is the vocabulary.
\end{itemize}

\textbf{Log-likelihood objective.}
\[
\max_\theta \sum_{(x,y)} \log p_\theta(y \mid x)
\]

For a single example,
\[
\log p_\theta(y \mid x) = \theta_y^\top \phi(x) - \log \sum_{y' \in V} \exp\!\left(\theta_{y'}^\top \phi(x)\right)
\]
Define the partition function
\[
Z(x) = \sum_{y' \in V} \exp\!\left(\theta_{y'}^\top \phi(x)\right)
\]
so that
\[
\log p_\theta(y \mid x) = \theta_y^\top \phi(x) - \log Z(x).
\]

\textbf{Gradient of the log-likelihood.} For a particular token \(y_0\),
\[
\nabla_{\theta_{y_0}} \log p_\theta(y \mid x)
= \left(\mathbf{1}[y=y_0] - p_\theta(y_0 \mid x)\right)\phi(x).
\]
This gives the key update intuition:
\begin{itemize}
    \item if \(y_0\) is the true next token, move \(\theta_{y_0}\) toward \(\phi(x)\);
    \item otherwise, move \(\theta_{y_0}\) away from \(\phi(x)\).
\end{itemize}

\textbf{Bag-of-words representations.}

Indicator bag-of-words maps a document to a \(V\)-dimensional vector whose entries indicate whether each token appeared.

Count bag-of-words uses token counts instead of binary indicators.

\textbf{TF-IDF.} The lecture motivates downweighting common function words.
\[
\operatorname{TF}(w,d) = \text{count of } w \text{ in document } d
\]
\[
\operatorname{IDF}(w) = \log \frac{N}{\operatorname{df}(w)}
\]
where \(N\) is the number of documents in the corpus and \(\operatorname{df}(w)\) is the number of documents containing \(w\). Then
\[
\operatorname{TFIDF}(w,d) = \operatorname{TF}(w,d)\operatorname{IDF}(w).
\]

\textbf{Skip-gram model.} Learn an input embedding \(w_x\) and an output embedding \(u_y\):
\[
p(y \mid x) = \frac{\exp\!\left(u_y^\top w_x\right)}
{\sum_{y' \in V} \exp\!\left(u_{y'}^\top w_x\right)}.
\]
This predicts a context word \(y\) from a center word \(x\).

\textbf{CBOW model.} Predict a target word from a set of context words:
\[
p(y \mid x) = \frac{\exp\!\left(u_y^\top \left(\frac{1}{|x|}\sum_{x_i \in x} w_{x_i}\right)\right)}
{\sum_{y' \in V} \exp\!\left(u_{y'}^\top \left(\frac{1}{|x|}\sum_{x_i \in x} w_{x_i}\right)\right)}.
\]
If we define
\[
W(x) = \frac{1}{|x|}\sum_{x_i \in x} w_{x_i},
\]
then
\[
\log p(y \mid x) = u_y^\top W(x) - \log \sum_{y' \in V} \exp\!\left(u_{y'}^\top W(x)\right).
\]

\textbf{CBOW gradient intuition.}
\[
\nabla_{u_{y_0}} \log p(y \mid x)
= \left(\mathbf{1}[y=y_0] - p(y_0 \mid x)\right) W(x).
\]
Thus co-occurring words and contexts are pulled together in embedding space, while non-co-occurring pairs are pushed apart.

\textbf{PMI connection.} For skip-gram and related embedding schemes, the learned embeddings approximately satisfy
\[
u_y^\top w_x \approx \log \frac{p(x,y)}{p(x)p(y)} - c
\]
for some constant \(c\). Thus dot products approximate \textbf{pointwise mutual information (PMI)}.

\subsection{Algorithms / Architectures}
\textbf{1. Build a fixed-feature log-linear LM}
\begin{itemize}
    \item Tokenize the corpus.
    \item Convert each training position into a supervised example \((x,y)\), where \(x\) is a context and \(y\) is the observed next token.
    \item Construct a fixed-dimensional context representation \(\phi(x)\).
    \item Score each candidate next token with \(\theta_y^\top \phi(x)\).
    \item Normalize scores with a softmax.
    \item Maximize log-likelihood with gradient-based optimization.
\end{itemize}

\textbf{2. Bag-of-words feature construction}
\begin{itemize}
    \item Start with token IDs.
    \item Build a vocabulary-sized vector.
    \item Fill it with either indicators or counts.
    \item Optionally apply TF-IDF weighting to reduce the impact of words that appear in nearly every document.
\end{itemize}

\textbf{3. Skip-gram training}
\begin{itemize}
    \item Choose a context window around each token.
    \item Treat the center word as input.
    \item Treat neighboring words as outputs.
    \item For each input-output pair, optimize the softmax probability of the output given the input.
    \item Learn input embeddings \(w\) and output embeddings \(u\).
\end{itemize}

\textbf{4. CBOW training}
\begin{itemize}
    \item Choose a context window around each target word.
    \item Average or sum the embeddings of the surrounding words.
    \item Use the resulting vector as the context representation.
    \item Predict the center word with a softmax over the vocabulary.
\end{itemize}

\textbf{5. Using learned embeddings downstream}
\begin{itemize}
    \item Train embeddings on large unlabeled text.
    \item Reuse them as input features for tasks such as sentiment classification, translation, or other NLP systems.
    \item This replaces hand-engineered sparse vectors with dense learned representations.
\end{itemize}

\subsection{Important Intuitions from the Professor}
\begin{itemize}
    \item \textbf{A strong language model must know much more than local syntax.} It implicitly needs grammar, factual knowledge, and common sense in order to assign high probability to real text.
    \item \textbf{Prompting is just conditional continuation.} This is the core modern perspective: many tasks become ``provide context, then continue it.''
    \item \textbf{The gradient has a geometric interpretation.} Training repeatedly moves parameter vectors toward contexts that support them and away from contexts that do not.
    \item \textbf{Bag-of-words is mostly about topical signal, not grammatical signal.} It can infer what a document is about but not what word should come next in a well-formed sentence.
    \item \textbf{Word2vec was historically important even though it is a poor generator.} Its importance was that it learned useful word representations, not that it generated fluent text.
    \item \textbf{The role of embeddings was originally transfer.} People first used these models to pretrain vectors that later improved other supervised NLP systems.
    \item \textbf{Finite windows are a real modeling choice.} Skip-gram and CBOW only learn relationships within a chosen context radius; anything outside that window is invisible.
    \item \textbf{Embedding geometry is not magic.} The lecture argues that analogies and manifolds arise because low-rank factorization of co-occurrence structure induces meaningful latent directions.
\end{itemize}

\subsection{Examples}
\begin{itemize}
    \item \textbf{Grammar example:}
    \[
    \text{``The cats under the sofa \_\_\_''}
    \]
    A good model should prefer ``purr'' over ``purrs.''
    \item \textbf{Factual knowledge example:}
    \[
    \text{``Daniel Akaka was born in \_\_\_''}
    \]
    The lecture contrasts ``Honolulu'' with ``Chicago.''
    \item \textbf{Common-sense example:}
    \[
    \text{``If you drop an egg it will \_\_\_''}
    \]
    A realistic model should favor ``break'' over ``bounce.''
    \item \textbf{Machine translation decomposition:} a target sentence can be chosen using a crude translation model plus a language model that prefers fluent target-language strings.
    \item \textbf{Baseball bag-of-words example:} counts of words such as ``Chicago,'' ``Cubs,'' and ``baseball'' reveal topic; TF-IDF suppresses overly common words such as ``the.''
    \item \textbf{Skip-gram / CBOW example:} from the Dickens sentence
    \[
    \text{``it is a far, far better rest that I go to, than I have ever known''}
    \]
    skip-gram predicts nearby words from a center word, while CBOW predicts the center word from neighboring words.
    \item \textbf{Embedding analogies:}
    \[
    \text{MIT} - \text{east\_coast} + \text{west\_coast} \approx \text{Caltech}
    \]
    \[
    \text{MIT} - \text{sciences} + \text{humanities} \approx \text{Harvard}
    \]
    \item \textbf{Country-city and geography structure:} pairs such as France/Paris and Italy/Rome occupy similar relative positions; nearby countries may also form smooth low-dimensional manifolds.
    \item \textbf{Color grounding example:} color words cluster in embedding space in ways that align with perceptual similarity, despite the model never seeing images.
\end{itemize}

\subsection{Common Pitfalls / Limitations}
\begin{itemize}
    \item \textbf{No word order in bag-of-words, skip-gram, or CBOW.} These models cannot represent syntax-sensitive phenomena that depend on position.
    \item \textbf{Sparse fixed features do not share input-side similarity.} ``good'' and ``great'' begin as unrelated coordinates unless the model learns some indirect output association.
    \item \textbf{Vocabulary-scale issues remain.} Sparse bag-of-words features still grow with vocabulary size.
    \item \textbf{Finite context windows lose long-range dependencies.} Words outside the training window do not affect the learned relation.
    \item \textbf{Poor open-ended generation.} Sampling from these models produces incoherent word salad because they lack order-sensitive compositional structure.
    \item \textbf{Embedding geometry is approximate and noisy.} Analogies and manifolds are trends, not exact symbolic facts.
    \item \textbf{Separate input and output embeddings add flexibility but also more parameters.} They are not always necessary.
\end{itemize}

\subsection{Connections to Other Lectures}
\begin{itemize}
    \item This lecture directly builds on Lecture 1's setup of language modeling as learning
    \[
    p_\theta(x_t \mid x_{1:t-1}).
    \]
    The new ingredient is a concrete parameterization.
    \item The fixed-feature log-linear model is the first answer to Lecture 1's open question: how should we parameterize next-token probabilities without explicit tables?
    \item The embedding discussion anticipates later neural models: dense continuous representations become the basic input objects for RNNs, attention, and Transformers.
    \item The lecture ends by highlighting the missing ingredient, \textbf{word order}, which motivates the next class and the broader move toward sequence models that truly model compositional structure over time.
    \item The PMI discussion foreshadows later themes in representation learning: much of modern NLP can be interpreted as learning compressed structure from co-occurrence statistics.
\end{itemize}

\subsection{Exam-Relevant Summary}
\begin{itemize}
    \item \textbf{Translation motivation:}
    \[
    p(y \mid x) \propto p(x \mid y)p(y)
    \]
    where \(p(y)\) is the language model.
    \item \textbf{Log-linear LM:}
    \[
    p_\theta(y \mid x) = \frac{\exp(\theta_y^\top \phi(x))}{\sum_{y'} \exp(\theta_{y'}^\top \phi(x))}
    \]
    \item \textbf{Log-likelihood form:}
    \[
    \log p_\theta(y \mid x) = \theta_y^\top \phi(x) - \log \sum_{y'} \exp(\theta_{y'}^\top \phi(x))
    \]
    \item \textbf{Gradient:}
    \[
    \nabla_{\theta_{y_0}} \log p_\theta(y \mid x)
    = \left(\mathbf{1}[y=y_0] - p_\theta(y_0 \mid x)\right)\phi(x)
    \]
    \item \textbf{Bag-of-words variants:} indicators, counts, TF-IDF.
    \item \textbf{TF-IDF intuition:} emphasize words frequent in one document but rare across the corpus.
    \item \textbf{Bag-of-words strengths:} topical information.
    \item \textbf{Bag-of-words weakness:} no word order, poor generation, sparse features.
    \item \textbf{Skip-gram:} predict context from a center word.
    \item \textbf{CBOW:} predict a target word from surrounding context.
    \item \textbf{Embeddings reduce parameters:} from roughly \(O(|V|^2)\) to \(O(|V|d)\) with embedding dimension \(d\).
    \item \textbf{PMI result:}
    \[
    u_y^\top w_x \approx \log \frac{p(x,y)}{p(x)p(y)} - c
    \]
    \item \textbf{Key limitation of all models in this lecture:} they do not model word order.
    \item \textbf{Historical importance:} these models were foundational because they learned reusable dense word representations that improved downstream NLP.
\end{itemize}

\end{document}
